% Part: computability
% Chapter: computability-theory
% Section: applications-fixed-point

\documentclass[../../../include/open-logic-section]{subfiles}

\begin{document}

\olfileid{cmp}{thy}{apf}
\olsection{Ngetrapake Teorema Titik Tetep}

Teorema titik tetep sejatine ngidini kita netepake
fungsi komputabel parsial adhedhasar indekse.
Tuladhane, kita bisa nemokake indeks $e$
saengga kanggo saben $y$,
\[
\cfind{e}(y) = e + y.
\]
Minangka tuladha liyane, bukti teorema titik tetep
bisa digunakake kanggo ngrancang program ing Java
utawa C++ kang nyithak program iku dhewe.

Elinga yen kanggo saben $e$, $W_e$ kita tetepake
minangka domain $\cfind{e}$, mula runtutan $W_0$,
$W_1$, $W_2$,~\dots ngenumerasi himpunan kang
bisa dienumerasi sacara komputabel. Sawetara
himpunan iki komputabel. Kita bisa takon apa ana
algoritma kang nampa nilai $x$ minangka input,
lan yen $W_x$ komputabel, mbalekake indeks
fungsi karakteristike. Wangsulane ``ora,''
algoritma kaya mangkono ora ana:

\begin{thm}
Ora ana fungsi komputabel parsial $f$ kang
nduwé sipat mangkene: saben $W_e$ komputabel,
$f(e)$ ditegesi lan $\cfind{f(e)}$ minangka
fungsi karakteristike.
\end{thm}

\begin{proof}
% OLPL-144: Sumber bukti nyebut f komputabel sanajan teorema nglimputi f komputabel parsial; nalika f(e) ora ditegesi, syarat teorema langsung gagal lan phi_{f(e)} mung diwaca kanthi kondisional.
Tetepna $f$ minangka fungsi komputabel parsial
sembarang; kita bakal mbangun $e$ saengga $W_e$
komputabel, nanging yen nilai f ing e ana,
$\cfind{f(e)}$ dudu fungsi karakteristike.
Kanthi nggunakake teorema titik tetep, kita
bisa nemokake indeks $e$ saengga
\[
\cfind{e}(y) \simeq
\begin{cases}
0 & \text{yen $y=0$ lan $\cfind{f(e)}(0) \fdefined = 0$} \\
\text{ora ditegesi} & \text{yen ora.}
\end{cases}
\]
Yaiku, $e$ diasilake kanthi ngetrapake
teorema titik tetep marang fungsi kang
ditetepake dening
\[
g(x,y) \simeq
\begin{cases}
0 & \text{yen $y=0$ lan $\cfind{f(x)}(0) \fdefined = 0$} \\
\text{ora ditegesi} & \text{yen ora.}
\end{cases}
\]
Sacara informal, kita bisa ndeleng yen $g$
komputabel parsial kanthi cara mangkene:
nalika nampa input $x$ lan $y$, algoritma
luwih dhisik mriksa apa $y$ padha
karo~$0$. Yen iya, algoritma ngitung
$f(x)$, banjur nganggo mesin universal
kanggo ngitung $\cfind{f(x)}(0)$.
Yen komputasi pungkasan iki mandheg
lan mbalekake~$0$, algoritma mbalekake~$0$;
yen ora, algoritma ora mandheg.

Nanging saiki gatekna yen $\cfind{f(e)}(0)$
ditegesi lan padha karo $0$, mula
$\cfind{e}(y)$ ditegesi persis nalika $y$
padha karo $0$, dadi $W_e =
\{ 0 \}$. Yen $\cfind{f(e)}(0)$ ora
ditegesi, kalebu yen nilai f ing e ora ana,
utawa ditegesi nanging ora padha karo
$0$, mula $W_e = \emptyset$. Ing loro
kahanan iki, yen nilai f ing e ora ana,
syarat teorema langsung gagal; yen ana,
$\cfind{f(e)}$ dudu fungsi karakteristik
saka~$W_e$, amarga menehi wangsulan
kang salah nalika nampa input $0$.
\end{proof}

\end{document}
